% =====================================================================
%  Energy Detective -- IEEE conference paper DRAFT (DIC4)
%  Group 20 | https://github.com/lun-09-15/AIoT_Final_Proj
%  Status: proposal stage. No system is implemented and no experiment
%  has been run. Every outcome is labeled "Expected"/"planned"; result
%  tables are placeholders (TBD). Do NOT fill in numbers before running.
% =====================================================================
\documentclass[conference]{IEEEtran}
\IEEEoverridecommandlockouts

\usepackage{cite}
\usepackage{amsmath,amssymb,amsfonts}
\usepackage{algorithm}
\usepackage{algpseudocode}
\usepackage{graphicx}
\usepackage{booktabs}
\usepackage{array}
\usepackage{multirow}
\usepackage{xcolor}
\usepackage{url}
\usepackage{listings}
\usepackage{tikz}
\usetikzlibrary{arrows.meta,positioning,fit,calc}

\newcommand{\yes}{\checkmark}
\newcommand{\part}{$\circ$}
\newcommand{\no}{--}
\newcommand{\tbd}{\textcolor{gray}{TBD}}

\lstset{basicstyle=\ttfamily\scriptsize,breaklines=true,frame=single,
        columns=fullflexible,xleftmargin=2pt,xrightmargin=2pt}

\begin{document}

\title{Energy Detective: An Evidence-Grounded LLM Agent
for Household Electricity Analysis}

% Specs: https://github.com/timeil0611/DIC4-Energy-Detective/tree/main/openspec
\author{
\IEEEauthorblockN{Sheng-Kang Tsai, Chun-Kai Tsai, Yu-Lun Wu, Ching-Lun Chiu}
\IEEEauthorblockA{\textit{Group 20, AIoT Course}\\
\textit{Department of Computer Science and Engineering}\\
\textit{National Chung Hsing University}\\
Taichung, Taiwan}
}

\maketitle

% ---------------------------------------------------------------------
\begin{abstract}
Large language model (LLM) agents make household electricity data
accessible through natural language, but a fluent answer can still quote
the wrong date, confuse power (kW) with energy (kWh), treat missing
records as zero, or present a statistical deviation as an appliance
fault. We present the design of \emph{Energy Detective}, a single-agent
system for five household tasks---consumption lookup, period comparison,
anomaly-candidate screening, next-hour forecasting, and evidence-based
recommendations---built on the principle that \emph{the LLM interprets
and coordinates, deterministic tools compute every number, and every
answer is checked against recorded evidence before it is shown}. The
design combines (i) guarded analytics tools with explicit unit,
completeness, and no-look-ahead contracts, (ii) an evidence registry
that links each number, chart, and recommendation to its data slice,
method, and source, and (iii) an external response validator with a
bounded budget of eight tool calls and one repair. Requirements are
captured as a verifiable OpenSpec contract of 11 capabilities with
SHALL/WHEN--THEN scenarios. We further specify a layered evaluation on
the UCI single-household dataset: a forecasting comparison against two
naive baselines, and a controlled ablation of the validator over 50
frozen tasks $\times$ 2 conditions $\times$ 3 runs. This paper reports
the design and evaluation protocol; implementation and experiments are
ongoing, and all outcomes are stated as expected rather than measured.
\end{abstract}

\begin{IEEEkeywords}
LLM agents, tool use, household energy, evidence traceability,
response verification, spec-driven development
\end{IEEEkeywords}

% =====================================================================
\section{Introduction}
% =====================================================================

Residential buildings account for roughly 70\% of building energy
demand worldwide, and household energy use differs widely across
climates, appliance stocks, and end uses~\cite{iea2025efficiency}.
Field experiments show that giving households information about their
consumption can change behavior: a randomized trial of real-time
feedback in Dutch households reported a 2.2\% reduction in electricity
use~\cite{boomsma2025feedback}, and a 24-week field experiment in
student housing found that both social-comparison nudges and
capability-building boosts improved consumption trends relative to a
control group~\cite{paunov2026nudges}. These results motivate better
ways for residents to \emph{understand} their data; they do not by
themselves show that a given analysis tool saves energy.

LLM agents are a natural interface for this need. Recent work equips
agents with analysis tools, device context, and knowledge retrieval for
building energy management~\cite{he2025contextaware}, or translates
natural-language requests into appliance schedules via native tool
calling~\cite{jonah2026agentic}. The same studies, however, report that
correctness varies by task and that an agent's explanation and its
executed actions must be checked separately.

\textbf{Three classes of error.} We focus on the following question:
\emph{how can a household-energy agent give answers whose numbers,
time ranges, and units are correct, whose analyses and recommendations
are traceable to actual data, and which disclose limits when the data
cannot support a conclusion?} We identify three classes of error that a
fluent answer can hide:
\begin{enumerate}
  \item \emph{Computational errors}: wrong date range, power reported
        as energy, partial observations summed as a complete total,
        or a percentage computed on a zero or incomplete baseline.
  \item \emph{Over-inference}: treating an anomaly candidate as waste
        or a fault, or attributing a sub-metered circuit's increase to a
        single appliance---which aggregate metering cannot
        support~\cite{hart1992nilm}.
  \item \emph{Category confusion}: presenting a forecast as an
        observation, or a hypothesis as an established cause.
\end{enumerate}

\textbf{Three lines of defense.} Energy Detective addresses these with
(1) \emph{tool computation}: all quantities come from guarded tools,
never from free-text arithmetic; (2) \emph{evidence recording}: every
tool result is stored with its data version, replay time, interval,
fields, quality status, method, and result ID; and (3) \emph{response
checking}: before output, an external validator compares the draft's
numbers, units, dates, directions, and citations with the recorded
evidence and allows at most one repair.

We do not claim that any of these elements is new on its own; our
contribution is their integration for household electricity analysis
and a protocol to test whether they help. The expected contributions
are:
\begin{itemize}
  \item A domain-grounded agent design for five household-electricity
        tasks with explicit tool contracts and inference boundaries
        (Sections~\ref{sec:problem}--\ref{sec:methods}).
  \item An evidence model that makes each answer auditable, and a
        bounded validate-and-repair loop (Section~\ref{sec:design}).
  \item A requirement specification written as a verifiable OpenSpec
        contract that links each claim to testable scenarios
        (Section~\ref{sec:openspec}).
  \item A layered evaluation protocol separating forecasting accuracy,
        task correctness, and evidence consistency
        (Section~\ref{sec:eval}).
\end{itemize}

The evaluation is organized around three research questions:
\textbf{RQ1}---Does a feature-based Random Forest forecast outperform
persistence and daily-seasonal-naive baselines (MAE, RMSE)?
\textbf{RQ2}---Does the agent complete the five task types with correct
numbers, units, and time ranges?
\textbf{RQ3}---With the LLM, tools, knowledge, and data held fixed,
does adding the external validator and one repair increase evidence
consistency and appropriate limitation handling, and at what cost in
latency, tool calls, and unnecessary refusals?

% =====================================================================
\section{Related Work}
% =====================================================================

We organize prior work into four research lines (Table~\ref{tab:matrix}).

\textbf{Energy feedback and behavior.} Real-time feedback reduced
household electricity use in a randomized trial~\cite{boomsma2025feedback},
and nudges and boosts both improved conservation trends in student
housing, without a significant difference between
them~\cite{paunov2026nudges}. At the system level, the IEA highlights
demand flexibility as a source of efficiency and grid
value~\cite{iea2025flexibility}. These studies evaluate behavioral
outcomes, not the correctness of an analytical answer; delivering an
answer, a resident acting on it, and measured savings must be validated
separately.

\textbf{Short-term load forecasting.} Boosting ensembles forecast
next-hour residential load with site-dependent
accuracy~\cite{muqtadir2025nowcasting}, and an industrial study of
household energy estimation chose an interpretable custom tree over a
more accurate Random Forest to balance error, transparency, and
business consistency~\cite{beraud2024explainable}. Forecast error,
however, does not measure whether a conversational recommendation is
supported by evidence.

\textbf{LLM agents for energy management.} He and
Jazizadeh~\cite{he2025contextaware} build context-aware agents that
combine data analysis, device interaction, and retrieval, reporting
77\% accuracy on energy-analysis queries and 49\% on cost management,
and noting the instability of generated analysis code. Jonah et
al.~\cite{jonah2026agentic} find native tool interfaces more reliable
than free-text parsing for scheduling, yet good average performance
does not guarantee constraint satisfaction in conflict scenarios.

\textbf{LLM tool use and reliability.} ReAct interleaves reasoning
with actions~\cite{yao2023react}, and Toolformer learns when to call
external tools~\cite{schick2023toolformer}. Hallucination remains a
core failure mode of generation~\cite{ji2023hallucination}.
Self-Refine improves outputs with self-generated
feedback~\cite{madaan2023selfrefine}, while CRITIC grounds
self-correction in tool-based verification~\cite{gou2024critic}. These
methods are domain-agnostic and do not encode energy-specific
contracts such as unit conversion, completeness, or sub-meter
attribution limits.

\textbf{Research gap.} Behavioral studies do not validate agent
analysis, forecasting studies do not evaluate answer correctness, and
energy agents show that computation, tool execution, and constraint
compliance can fail. Generic verification methods lack domain rules.
We therefore study, in a single-household replay setting, whether
tool-computed numbers, an evidence registry, and a domain-aware
external validator make answers traceable and appropriately bounded.

\begin{table}[t]
\centering
\caption{Feature comparison across the four research lines
(\yes: addressed; \part: partially; \no: not a stated focus).
The last row describes \emph{planned} capabilities.}
\label{tab:matrix}
\setlength{\tabcolsep}{3pt}
\footnotesize
\begin{tabular}{@{}lcccccc@{}}
\toprule
Work & NL & Tool & Evid. & Ext. & Limit & House-\\
     & interf. & comp. & trace & verif. & discl. & hold \\
\midrule
Feedback RCT~\cite{boomsma2025feedback}   & \no & \no & \no & \no & \no & \yes \\
Nudge/boost~\cite{paunov2026nudges}       & \no & \no & \no & \no & \no & \yes \\
Boosting STLF~\cite{muqtadir2025nowcasting} & \no & \no & \no & \no & \no & \yes \\
Explainable est.~\cite{beraud2024explainable} & \no & \no & \part & \no & \no & \yes \\
Context-aware agent~\cite{he2025contextaware} & \yes & \yes & \no & \no & \no & \yes \\
Agentic HEMS~\cite{jonah2026agentic}      & \yes & \yes & \no & \part & \no & \yes \\
ReAct~\cite{yao2023react}                 & \yes & \yes & \part & \no & \no & \no \\
Self-Refine~\cite{madaan2023selfrefine}   & \yes & \no & \no & \part & \no & \no \\
CRITIC~\cite{gou2024critic}               & \yes & \yes & \part & \yes & \no & \no \\
\midrule
\textbf{Ours (planned)}                   & \yes & \yes & \yes & \yes & \yes & \yes \\
\bottomrule
\end{tabular}
\\[2pt]
\raggedright\scriptsize
NL interf.: natural-language interface; Tool comp.: numbers computed by
tools; Evid. trace: per-claim evidence links; Ext. verif.: output
checked by an external verifier; Limit discl.: explicit handling of
insufficient data. % TODO(team): re-check each cell against the papers.
\end{table}

% =====================================================================
\section{Problem Formulation and Design Principles}
\label{sec:problem}
% =====================================================================

\textbf{Data.} We use the UCI Individual Household Electric Power
Consumption dataset~\cite{hebrail2006uci}: minute-level records from a
single household in France (2006--2010), including global active power
(kW, one-minute average) and three sub-metered circuits (Wh per
minute). Each sub-meter aggregates several appliances.

\textbf{Historical replay.} The agent operates at a fixed replay
boundary $h$ (an hour boundary on the original clock). Only records
with timestamp $t<h$ are visible; relative terms such as ``today'' and
``yesterday'' are interpreted with respect to $h$. The forecast target
is the energy $y_h$ in $[h, h+1)$.

\textbf{Energy and completeness.} For global active power $P_i$ (kW)
and sub-meter energy $e_i$ (Wh) over minutes $i$ in an interval,
\begin{equation}
E_{\text{global}} = \sum_i \frac{P_i}{60}\ \text{kWh},\qquad
E_{\text{sub}} = \frac{1}{1000}\sum_i e_i\ \text{kWh}.
\end{equation}
An hour is \emph{complete} only when all 60 distinct minutes hold valid,
non-conflicting values. Missing values are never treated as zero and are
not imputed; an incomplete interval yields a \texttt{PARTIAL\_OBSERVED}
quantity with its coverage ratio. A percentage change is reported only
for a complete, non-zero baseline.

\textbf{Five core principles.} The design follows five principles
derived from our domain analysis:
(P1)~\emph{measurement definition}---power and energy are distinct;
(P2)~\emph{data granularity}---sub-meters cannot identify individual
appliances~\cite{hart1992nilm};
(P3)~\emph{anomaly $\neq$ cause}---a deviation from history is a
candidate for inspection, not a diagnosis;
(P4)~\emph{no look-ahead}---data views, features, and model training
cut-offs respect $h$;
(P5)~\emph{effect requires separate validation}---claims of energy,
cost, or carbon savings require intervention studies.

% =====================================================================
\section{System Design}
\label{sec:design}
% =====================================================================

\begin{figure}[t]
\centering
\resizebox{\columnwidth}{!}{%
\begin{tikzpicture}[
  box/.style={draw,rounded corners=2pt,align=center,font=\scriptsize,
              minimum height=7mm,inner sep=2pt,text width=21mm},
  arr/.style={-{Stealth[length=1.6mm]},thin},
  lab/.style={font=\tiny,fill=white,inner sep=1pt}]
\node[box] (q) at (0,0)        {\textcircled{\tiny 1} Query \&\\replay setting};
\node[box] (a) at (2.75,0)     {\textcircled{\tiny 2} Planning \&\\coordination agent};
\node[box] (u) at (5.5,0)      {\textcircled{\tiny 8} Streamlit UI:\\answer, chart, limits};
\node[box] (d) at (0,-1.35)    {\textcircled{\tiny 3} Data access\\\& quality check};
\node[box] (t) at (2.75,-1.35) {\textcircled{\tiny 4} Analytics \&\\forecast tools};
\node[box] (k) at (5.5,-1.35)  {\textcircled{\tiny 5} Energy\\knowledge base};
\node[box] (e) at (2.75,-2.6)  {\textcircled{\tiny 6} Evidence\\registry};
\node[box] (v) at (2.75,-3.85) {\textcircled{\tiny 7} Response \&\\evidence validator};
\draw[arr] (q) -- (a);
\draw[arr,dashed] (a) -- node[lab,above=2pt]{clarify} (u);
\draw[arr] (a) -- (t);
\draw[arr] (a.south west) -- (d.north east);
\draw[arr] (a.south east) -- (k.north west);
\draw[arr] (d) -- (t);
\draw[arr] (t) -- (e);
\draw[arr] (d.south) |- (e.west);
\draw[arr] (k.south) |- (e.east);
\draw[arr] (e) -- (v);
\draw[arr] (v.east) -- (6.85,-3.85) |- node[lab,pos=0.25,right=1pt,align=left]{pass /\\partial /\\limited} (u.east);
\draw[arr] (v.west) -- (-1.4,-3.85) |- node[lab,pos=0.25,left=1pt,align=right]{repair\\($\le$1)} (2.75,0.75) -- (a.north);
\end{tikzpicture}}
\caption{Architecture of Energy Detective (eight modules). Numbers are
computed only by \textcircled{\tiny 3}/\textcircled{\tiny 4}; every
result is logged in \textcircled{\tiny 6} and checked by
\textcircled{\tiny 7} before display.}
\label{fig:arch}
\end{figure}

Fig.~\ref{fig:arch} shows the eight modules. A single cloud LLM with
structured tool calling (\textcircled{\tiny 2}) interprets the question,
asks for clarification when the request is ambiguous, selects tools,
and drafts the answer; it never computes quantities itself.

\textbf{Tool contracts.} Table~\ref{tab:tools} lists the tools. Data
checks are enforced in the tool layer rather than relying on the LLM:
the agent cannot change $h$ or the read-only view, and every call
returns a status (\emph{success}, \emph{partial}, \emph{insufficient},
\emph{rejected}, or \emph{failed}) that is logged as evidence.

\begin{table}[t]
\centering
\caption{Guarded tool contracts (design names).}
\label{tab:tools}
\footnotesize
\begin{tabular}{@{}>{\raggedright\arraybackslash\ttfamily\scriptsize}p{0.40\columnwidth}p{0.54\columnwidth}@{}}
\toprule
\normalfont\footnotesize Tool & Required output / precondition \\
\midrule
ReadHistoricalData & Records with $t<h$ only, raw units, quality status, data version \\
SummarizeEnergy & kWh, coverage, completeness label, per-hour results \\
ComparePeriods & Absolute difference; percentage only for a complete, non-zero baseline \\
FlagAnomalyCandidates & Reference size, median, MAD, candidate flag, limitations \\
ForecastNextHour & Model and feature cut-offs, target $[h,h{+}1)$, method, estimate \\
RetrieveDomainKnowledge & Knowledge ID, rule, applicability limits, source; explicit ``no match'' \\
\bottomrule
\end{tabular}
\end{table}

\textbf{Evidence registry.} Each record stores the data version, replay
time, interval, fields, quality status, tool or model version,
parameters, result value and unit, and a result ID. Draft sentences and
charts must reference existing result IDs; the binder never attaches a
plausible-looking citation to an unsupported claim. The model's
internal reasoning is not treated as evidence.

\textbf{Response validator.} The validator checks required fields,
dates, units, numbers (tolerance 0.005 for two-decimal rounding plus a
small floating-point epsilon), direction of change, chart IDs, and
explicit boundary violations (e.g., attributing a sub-meter increase to
one appliance). Free-text reasoning that rules cannot capture is left to
human rating (Section~\ref{sec:eval}). Answers are labeled
\texttt{ANSWER}, \texttt{CLARIFY}, \texttt{PARTIAL}, or
\texttt{LIMITED}; these are statuses, not confidence scores.

\textbf{Bounded control loop.} Algorithm~\ref{alg:agent} summarizes the
loop. Each task allows at most eight agent-initiated tool calls (failed
calls count) and one repair; neither budget resets. When the budget is
exhausted or validation fails twice, a fixed-structure terminal reply
reports only confirmed evidence and the reason for stopping, so an
unchecked draft is never shown.

\begin{algorithm}[t]
\caption{Energy Detective agent loop}
\label{alg:agent}
\footnotesize
\begin{algorithmic}[1]
\Require query $Q$, replay time $h$, read-only view $V$
\Ensure evidence-linked response $R$
\State $C \gets \Call{FreezeContext}{h, V, \text{versions}}$
\State $E \gets \emptyset$;\ $calls \gets 0$;\ $repairs \gets 0$
\Loop
  \State $A \gets \Call{PlanNextAction}{Q, C, E}$
  \If{$A.\mathit{type}=\textsc{Clarify}$}
     \State \Return \Call{RenderClarification}{A}
  \ElsIf{$A.\mathit{type}=\textsc{Tool}$}
     \If{$calls \ge 8$} \State \Return \Call{SafeTerminal}{E, \text{budget}} \EndIf
     \State $calls \gets calls+1$
     \State $T \gets \Call{ExecuteGuardedTool}{A, C}$;\ $E \gets E \cup \{T\}$
  \ElsIf{$A.\mathit{type}=\textsc{Draft}$}
     \State $D \gets \Call{BindDraftToEvidence}{A.\mathit{draft}, E}$
     \State $r \gets \Call{ValidateResponse}{D, E, C}$
     \If{$r.\mathit{pass}$} \State \Return \Call{Render}{D, E} \EndIf
     \If{$repairs \ge 1$} \State \Return \Call{SafeTerminal}{E, \text{invalid}} \EndIf
     \State $repairs \gets repairs+1$;\ $Q \gets Q \oplus r.\mathit{issues}$
  \Else
     \State \Return \Call{SafeTerminal}{E, \text{error}}
  \EndIf
\EndLoop
\end{algorithmic}
\end{algorithm}

% =====================================================================
\section{Analytics and Domain Knowledge}
\label{sec:methods}
% =====================================================================

\textbf{Next-hour forecasting.} We compare a Random Forest
regressor~\cite{pedregosa2011sklearn} (initial setting: 200 trees,
max depth 12, min.\ 5 samples per leaf) against \emph{persistence}
($\hat y_h = y_{h-1}$) and \emph{daily seasonal naive}
($\hat y_h = y_{h-24}$). Features are lags $y_{h-1}, y_{h-2},
y_{h-24}, y_{h-168}$; the mean and standard deviation of the last 3 and
24 complete hours; and calendar features known at $h$. Splits are
chronological (70/15/15) on the full time grid; the test period uses
rolling one-step replay without re-tuning, and the model's training
cut-off may never exceed $h$. These hyperparameters are a starting
point, not tuned results.

\textbf{Anomaly candidates.} For a target hour value $x$, the reference
set contains complete readings at the same hour and day type
(weekday/weekend) in the preceding 56 days, requiring at least 14
samples. With median $m$ and $\mathrm{MAD}=\mathrm{median}(|x_i-m|)$,
the modified z-score~\cite{iglewicz1993outliers,nist_outliers}
\begin{equation}
M = \frac{0.6745\,(x-m)}{\mathrm{MAD}}
\end{equation}
flags a candidate when $|M|>3.5$. If the reference is too small or
$\mathrm{MAD}=0$, the tool returns a limitation instead of a score.
Because the dataset has no fault or waste labels, we do not report
detection accuracy; we instead report flag frequencies at thresholds
2.5/3.5/4.5 as a sensitivity analysis.

\textbf{Domain knowledge.} A small curated base of about ten entries
(K01--K10) covers units, sub-meter granularity, missing data, comparison
baselines, zero denominators, anomaly candidates, forecast limits, the
difference between energy saving, cost saving, and peak shaving,
behavioral effects, and generalization. Each entry records a
definition, applicability conditions, what \emph{cannot} be inferred,
and its source. Retrieval is concept/keyword based; an empty match is
reported rather than fabricated.

% =====================================================================
\section{Spec-Driven Development with OpenSpec}
\label{sec:openspec}
% =====================================================================

To turn design principles into verifiable requirements, we adopted
OpenSpec~\cite{openspec}, a spec-driven development framework for AI
coding assistants. The project configuration records the research
context and writing rules (e.g., all outcomes must be labeled
\emph{Expected}; no fabricated results). The design document was
converted into one OpenSpec change consisting of a proposal,
\textbf{11 capability specifications}, \textbf{9 design decisions}, and
\textbf{32 implementation tasks}. Every requirement is phrased with
SHALL and tested by WHEN/THEN scenarios, and the change passes OpenSpec's
strict validation. The specification is public at
\url{https://github.com/timeil0611/DIC4-Energy-Detective}. A traceability table maps each claim in this paper
(principles P1--P5, the three trade-offs in Section~\ref{sec:disc}, and
RQ1--RQ3) to at least one capability. An abridged excerpt from the \texttt{agent-orchestration} and \texttt{data-quality} specifications:

\begin{lstlisting}
### Requirement: Bounded tool budget
The agent SHALL execute at most 8 agent-initiated
tool calls per task, counting failed and rejected
calls, and the budget SHALL NOT reset after a repair.

#### Scenario: Budget exhausted
- WHEN a task has used 8 tool calls and the agent
  requests another tool
- THEN the system returns a LIMITED reply built only
  from recorded evidence

### Requirement: Complete hour definition
#### Scenario: 59 valid minutes
- WHEN an hour has 59 valid minutes for the field
- THEN the hour is labeled PARTIAL_OBSERVED with
  coverage 59/60
\end{lstlisting}

These scenarios double as acceptance tests: during implementation and
evaluation, each behavior can be checked directly against the
specification rather than against prose.

% =====================================================================
\section{Evaluation Protocol}
\label{sec:eval}
% =====================================================================

\textbf{RQ1: forecasting.} The three methods are compared on common
test hours with complete targets, using MAE and RMSE (kWh), together
with each method's availability rate and sample count. Excluded samples
are reported with reasons.

\textbf{RQ2/RQ3: agent tasks.} A pilot set of 10 tasks is used only to
tune tool descriptions. The frozen formal set has 50 tasks: 10 per task
type, of which 6 are regular and 4 are boundary cases (missing data,
ambiguous time, unidentifiable appliance). Each task fixes its replay
time, data view, and expected answer or expected clarification. Two
conditions share the same LLM version, tools, knowledge, data, and
prompts: \textbf{A}---no external validation; \textbf{B}---validator
with at most one repair, under the same tool budget. Each task is run
three times per condition ($50\times2\times3=300$ task executions; a
task may involve several model calls).

\textbf{Metrics.} Table~\ref{tab:metrics} lists the metrics. To avoid
rewarding evasive answers, evidence consistency is always reported
together with coverage, task completion, and the unnecessary-refusal
rate on regular tasks. Two team members independently rate
recommendations on evidence support, inference boundary, and
actionability (0--2 each) and report initial agreement; LLM
self-assessment is not used as the sole judge.

\begin{table}[t]
\centering
\caption{Planned metrics and result placeholders (no experiment has
been run; all cells are TBD).}
\label{tab:metrics}
\footnotesize
\begin{tabular}{@{}lcc@{}}
\toprule
Metric & Cond.\ A & Cond.\ B \\
\midrule
Task completion rate (50 tasks)          & \tbd & \tbd \\
Number / unit / time correctness         & \tbd & \tbd \\
Evidence consistency                     & \tbd & \tbd \\
Correct limitation handling (20 boundary)& \tbd & \tbd \\
Unnecessary refusal (30 regular)         & \tbd & \tbd \\
Recommendation score (0--6, human)       & \tbd & \tbd \\
Latency / tool calls / repairs per task  & \tbd & \tbd \\
\midrule
\multicolumn{3}{@{}l}{RQ1: MAE / RMSE (kWh) of RF, persistence,
seasonal naive: \tbd} \\
\bottomrule
\end{tabular}
\end{table}

We pre-register no expected improvement. If condition B does not
improve, is slower, or refuses more often, this will be reported and
analyzed, and all runs---not only the best---will be reported.

% =====================================================================
\section{Discussion: Trade-offs and Limitations}
\label{sec:disc}
% =====================================================================

\textbf{Three design trade-offs.}
(i)~\emph{Social comparison vs.\ concrete guidance}: comparative
feedback can motivate, but a single-household dataset offers no peer
group, so we provide concrete, evidence-linked inspection suggestions
instead~\cite{paunov2026nudges}.
(ii)~\emph{Forecast accuracy vs.\ transparency}: a Random Forest may be
more accurate than naive baselines but less transparent; industry
practice sometimes favors interpretable models~\cite{beraud2024explainable},
so baselines are always reported alongside the model.
(iii)~\emph{LLM autonomy vs.\ deterministic checking}: more autonomy
can improve flexibility but makes errors harder to audit; we constrain
computation to tools and add rule-based checks, accepting extra latency
that RQ3 will quantify.

\textbf{Limitations.} A single, older French household does not
represent other households or Taiwan; the data lack appliance-level
activity, weather, and applicable tariffs; there are no ground-truth
fault labels, so anomalies remain candidates; rule-based checks cannot
judge all free-text reasoning; thresholds and budgets are not yet
validated on this data; and cloud LLM versions may affect
reproducibility. Without adoption and before/after measurements, we
make no claim of energy, cost, or carbon savings.

% =====================================================================
\section{Conclusion}
% =====================================================================

We presented the design of Energy Detective, an LLM agent for
household electricity analysis in which the LLM coordinates,
deterministic tools compute, and an external validator checks every
answer against recorded evidence. The requirements are captured as an
OpenSpec contract with testable scenarios, and a controlled evaluation
protocol is specified. Next steps are to implement the data and
analytics tools, select and pin a cloud LLM, build the validator and
Streamlit interface, run the pilot, freeze the formal task set, and
report results for RQ1--RQ3.

% ---------------------------------------------------------------------
\section*{Acknowledgment}
% TODO: course / instructor acknowledgment if required.
This draft was prepared with AI assistance (MCP-connected coding
assistant) for drafting and specification; all design decisions and
sources were reviewed by the authors.

\bibliographystyle{IEEEtran}
\bibliography{refs}

\end{document}
